\begin{frame}{같은 무작위 걷기, 다른 질문}
  \small
  \begin{columns}[T]
    \begin{column}{0.5\textwidth}
      \textbf{\color{ksablue}13.2 Node2Vec}
      \begin{itemize}
        \item ``이 걷기가 만들어내는 \textbf{노드의
          순서(시퀀스)}는 어떤 구조를 담고 있는가?''
        \item 그 시퀀스를 ``문장'' 취급해 word2vec의
          skip-gram을 학습 $\to$ \textbf{노드 임베딩}.
      \end{itemize}
    \end{column}
    \begin{column}{0.5\textwidth}
      \textbf{\color{ksablue}13.3 PageRank}
      \begin{itemize}
        \item ``이 걷기를 영원히 돌리면 \textbf{각 노드에 머무는
          확률} $\pi^*$는 얼마인가?''
        \item 그 정상분포 자체를 노드의 ``\textbf{중요도}''로 읽는다.
      \end{itemize}
    \end{column}
  \end{columns}
  \vskip0.6em
  \begin{block}{한 줄}
    같은 재료(그래프, 무작위 걷기)에서 하나는 \textbf{시퀀스}를 뽑아내고,
    하나는 \textbf{정상분포}를 뽑아낸다 --- 벡터와 스칼라,
    두 가지 다른 표현.
  \end{block}
\end{frame}
